\iffalse
\begin{corollary}
Кількість нільпотентних клітин Жордана задається числом $\dim (\ker B)$. Тобто в даному випадку, всього $p$ клітин, бо $\dim (\ker B) = p$.
\end{corollary}
\fi

%GB approach
\iffalse
\subsection{Приєднаний власний вектор (старий)}
\begin{remark}
Надалі ми будемо враховувати комплексні корені з характеристичного полінома, щоб ми мали змогу завжди знайти власні числа.
\end{remark}

\iffalse
Розглянемо особливий випадок, коли $A: L \to L$ має єдиний (взагалі єдиний л.н.з.) власний вектор $f$, що відповідає власному числу $\lambda$. При цьому $\dim L = n$. Це один з випадків, де оператор $A$ не може бути діагоналізованим.\\
За вимогою, $\dim (\ker(A - \lambda I)) = \dim L_{\lambda} = 1$. Отже, за рівністю про зв'язок з ядром та образом, $\dim (\Im (A - \lambda I)) = n-1$.\\
Задамо таку множину $M=\ker(A-\lambda I) \cap \Im(A-\lambda I)$. Оскільки $M \subset \ker (A-\lambda I)$, то за одною лемою, $\dim M \leq \dim(\ker (A-\lambda I)) = 1$.
Тоді звідси або $\dim M = 0$, або $\dim M = 1$.
\bigskip \\
Перед цим доведемо, що $\ker(A - \lambda I), \Im(A-\lambda I)$ - інваріантні підпростори для оператора $A$. За вимогою, $\ker(A-\lambda I) = L_{\lambda} = \linspan\{f\}$.\\
Тоді $\forall g \in \linspan\{f\}: g = \alpha f \implies Ag = A(\alpha f) = \alpha Af = (\alpha \lambda) f \implies Ag \in \linspan\{f\}$. Або теж саме, що $\forall g \in (\ker(A - \lambda I)) \implies Ag \in  \ker(A - \lambda I))$.
\bigskip \\
Тоді $\forall y \in \Im(A-\lambda I): \exists x \in L: y = (A-\lambda I)x \implies \\ Ay = A(A-\lambda I)x = (A^2-\lambda A)x = (A-\lambda I)(Ax) \in \Im(A-\lambda I)$.\\
Отже, $\ker(A-\lambda I)$ та $\Im(A-\lambda I)$ - два інваріантних підпростори.
\bigskip \\
Розглянемо випадок $\dim M = 0$. Звідси маємо $M = \ker(A-\lambda I) \cap \Im(A-\lambda I) = \{0\}$.\\
Розглянемо $A |_{\ker(A-\lambda I)}: \ker(A-\lambda I) \to \ker(A-\lambda I)$ - звужений оператор. У нього є власний вектор $f$.\\
Також розглянемо $A |_{\Im(A-\lambda I)}: \Im(A-\lambda I) \to \Im(A-\lambda I)$ - звужений оператор. У цього оператора є власне число $\mu$ та власний вектор $g$ (згідно з нашим зауваженням).\\
$Ag = A |_{\Im(A-\lambda I)} g = \mu g$\\
Таким чином, $f \in \ker(A-\lambda I)$ та $g \in \Im(A-\lambda I)$. Оскільки $\ker \{ A-\lambda I \} \cap \Im \{ A - \lambda I \} = \{ 0 \}$, то вони утворюють пряму суму, а водночас $\{f,g\}$ - л.н.з. і є власними векторами для $A$ - суперечність. Отже, $\dim M \neq 0$.
\bigskip \\
Залишається єдиний випадок - це $\dim M = 1$.\\
$\begin{cases}
M \subset \ker(A-\lambda I)\\
\dim{(\ker(A-\lambda I))} = 1
\end{cases}
\implies M = \ker(A-\lambda I)$
\bigskip \\
Тоді як можна побачити, $\ker(A-\lambda I) = \ker(A-\lambda I) \cap \Im (A - \lambda I)$, тобто $\ker(A-\lambda I) \subset \Im(A-\lambda I)$\\
$f \in \ker{A- \lambda I} \implies f \in \Im(A - \lambda I)$\\
Тоді $\exists h \in L: f = (A-\lambda I)h$. Ми прийшли до нового означення:

\begin{definition}
Задано $A: L \to L$ - лінійний оператор, в якому $f$ - власний вектор з власним числом $\lambda$.\\
Елемент $h \in L$ називається \textbf{приєднаним вектором} до власного вектора $f$ (висоти 1), якщо
\begin{align*}
(A-\lambda I)h = f
\end{align*}
\end{definition}

\begin{lemma}
Задано $A: L \to L$ - лінійний оператор. Він має єдиний власний вектор $f$ та $\dim L > 1$. Тоді для $f$ існує приєднаний власний вектор $h$.\\
\textit{Доведення було перед означенням приєданого вектора.}
\end{lemma}

\begin{definition}
Задано $A: L \to L$ - лінійний оператор, в якому $f$ - власний вектор з власним числом $\lambda$.\\
Елемент $h^{(k)} \in L$ називається \textbf{приєднаним вектором} до власного вектора $f$ \textbf{висоти $k$}, якщо
\begin{align*}
(A-\lambda I)h^{(k)} = h^{(k-1)}
\end{align*}
\end{definition}

\begin{remark}
Буду позначати $h^{(1)} = h$.
\end{remark}

\begin{remark}
Рівняння для приєднаного висоти $k$ можна записати наступним чином:\\
$(A-\lambda I)f = 0$\\
$(A-\lambda I)h^{(1)} = f \Rightarrow (A-\lambda I)^2 h^{(1)} = 0$\\
$(A-\lambda I)h^{(2)} = h \Rightarrow (A-\lambda I)^3 h^{(2)} = 0$\\
$\vdots$\\
$(A-\lambda I)^{k+1}h^{(k)} = 0$\\
Остаточно отримаємо таку форму:
\begin{align*}
(A-\lambda I)^{k+1}h^{(k)} = 0
\end{align*}
Взагалі-то кажучи, можна це проілюструвати ось так
\begin{figure}[H]
\centering
\begin{tikzcd}
0 & f \arrow{l}{A-\lambda I} & h^{(1)} \arrow{l}{A-\lambda I} & h^{(2)} \arrow{l}{A-\lambda I} & \dots \arrow{l}{A-\lambda I} & h^{(k)} \arrow{l}{A-\lambda I}
\end{tikzcd}
\end{figure}
\end{remark}

\begin{theorem}
Задано $A: L \to L$ - лінійний оператор, в якому $f$ - власний вектор з власним числом $\lambda$. Відомо, що $\{f,h,h^{(2)},\dots,h^{(k)}\}$ - ланцюг власного та приєднаних до нього векторів.\\
Тоді $\{f,h,h^{(2)},\dots,h^{(k)}\}$ - л.н.з.
\end{theorem}

\begin{proof}
$\alpha_0 f + \alpha_1 h^{(1)} + \dots + \alpha_k h^{(k)} = 0$.\\
Подіємо оператором $(A-\lambda I)$ покроково $k$ разів:\\
$(A-\lambda I)(\alpha_0 f + \alpha_1 h^{(1)} + \dots + \alpha_k h^{(k)}) = 0$\\
$\alpha_1 f + \alpha_2 h^{(1)} + \dots + \alpha_k h^{(k-1)} = 0$\\
$\alpha_2 f + \dots + \alpha_k h^{(k-2)} = 0$\\
$\dots$\\
$\alpha_{k-1} f + \alpha_k h^{(1)} = 0$\\
$\alpha_k f = 0$\\
$\Rightarrow \alpha_k = 0 \Rightarrow \alpha_{k-1} = 0 \Rightarrow \dots \Rightarrow \alpha_2 = 0 \Rightarrow \alpha_1 = 0 \Rightarrow \alpha_0 = 0$.\\
Отже, система $\{f,h,h^{(2)},\dots,h^{(k)}\}$ - л.н.з.
\end{proof}

\begin{theorem}
Задано $A: L \to L$ - лінійний оператор, в якому $f$ - єдиний власний вектор з власним числом $\lambda$ та нехай $\dim L = n$. Відомо, що $\{f,h,h^{(2)},\dots,h^{(k)}\}$ - ланцюг власного та приєднаних до нього векторів.
Тоді $\{f,h,\dots,h^{(n-1)}\}$ - базис в просторі $L$.
\end{theorem}

\begin{proof}
Лінійна незалежність вже є. Але єдине, що залишається довести, - це те, що існує ланцюг саме довжини $n$\\
Ми вже в курсі, що $\dim(\ker(A-\lambda I)) = 1$ та $\Im(A-\lambda I) \overset{\textrm{позн.}}{=} L_1$ - інваріантний підпростір для $A$, $\dim L_1 = n-1$\\
Більш того, $\ker(A-\lambda I) \subset \Im(A-\lambda I)$
\bigskip \\
Розглянемо оператор $A_1 = A |_{L_1}: L_1 \to L_1$\\
У нього єдиний л.н.з. власний вектор $f$, оскільки $\ker(A-\lambda I) \subset L_1 \subset L$\\
Тоді $\ker(A_1-\lambda I) = \ker(A-\lambda I) = \linspan\{f\}$\\
Тому $\dim(\ker(A_1 - \lambda I)) = 1 \Rightarrow \dim (\Im(A_1-\lambda I)) = (n-1)-1=n-2$\\
$L_2 = \Im(A_1-\lambda I)$ - інваріантний підпростір для $A_1$, а отже, й для $A$. Доводиться аналогічним чином як на початку цього пункту\\
Зауважимо, що $\forall y \in L_2: \exists z \in L_1: \\ y = (A_1 - \lambda I)z = (A-\lambda I)z \boxed{=}$\\
$z \in L_1 = \Im(A-\lambda I) \Rightarrow \exists x \in L: z = (A-\lambda I)x$\\
$\boxed{=} (A-\lambda I)^2 x$\\
Отримали, що $\Im(A_1 - \lambda I) = \Im(A-\lambda I)^2 = L_2$, $\dim L_2 = n-2$
\bigskip \\
Розглянемо оператор $A_2 = A |_{L_2}: L_2 \to L_2$\\
У нього єдиний л.н.з. власний вектор $f$ за аналогічними міркуваннями\\
Тоді $\ker(A_2-\lambda I) = \ker(A-\lambda I) = \linspan\{f\}$\\
Тому $\dim(\ker(A_2 - \lambda I)) = 1 \Rightarrow \dim (\Im(A_2-\lambda I)) = (n-2)-1=n-3$\\
$L_3 = \Im(A_2-\lambda I)$ - інваріантний підпростір для $A_2$, а отже, й для $A$\\
Отримаємо, що $\dim (A_2 - \lambda I) = \dim (A-\lambda I)^3$\\
І знову теж саме...
\bigskip \\
Тоді остаточно, $L \supset \underset{=\Im(A-\lambda I)}{L_1} \supset \underset{=\Im(A-\lambda I)^2}{L_2} \supset \dots \supset \underset{=\Im(A-\lambda I)^{n-1}}{L_{n-1}} \supset \{0\}$\\
Зауважимо,\\
$\begin{cases}
\dim (\Im(A-\lambda I)^{n-1}) = 1\\
\dim (\ker (A-\lambda I)) = 1\\
\ker (A-\lambda I) \subset \Im(A-\lambda I)^{n-1}
\end{cases} \Rightarrow \Im(A-\lambda I)^{n-1} = \ker(A-\lambda I)
$\\
Знайдемо ланцюг власного та приєднаного векторів довжини $n$\\
$f$ - власний\\
$f \in L_{n-1} \Rightarrow f \in L_1 \Rightarrow \exists h^{(1)} \in L_1 \Rightarrow h^{(1)} \in L: (A-\lambda I)h = f$\\
$f \in L_{n-1} \Rightarrow f \in L_2 \Rightarrow \exists h^{(2)} \in L_2 \Rightarrow h^{(2)} \in L: (A-\lambda I)^2 h^{(2)} = f$\\
$\dots$\\
$f \in L_{n-1} \Rightarrow \exists h^{(n-1)} \in L_{n-1} \Rightarrow h^{(n-1)} \in L: (A-\lambda I)^{(n-1)}h^{(n-1)} = f$
\end{proof}

Тоді можемо отримати матрицю оператора $A: L \to L$ в базисі $\{f,h,h^{(2)},\dots,h^{(k)}\}$\\
$(A-\lambda I)f = 0$\\
$(A-\lambda I)h = f = f + 0h + \dots + 0h^{(k)}$\\
$(A-\lambda I)h^{(2)} = h = 0f + h + \dots + 0h^{(k)}$\\
$\dots$\\
$(A-\lambda I)h^{(k)} = h^{(k-1)} = 0f + 0h + \dots + h^{(k-1)} + 0h^{(k)}$\\
$(\mathbb{A}- \lambda \mathbb{I}) = \begin{pmatrix}
0 & 1 & 0 & 0 & \dots & 0 \\
0 & 0 & 1 & 0 & \dots & 0 \\
0 & 0 & 0 & 1 & \dots & 0 \\
\vdots & \vdots & \vdots & \vdots & \ddots & \vdots \\
0 & 0 & 0 & 0 & \dots & 1 \\
0 & 0 & 0 & 0 & \dots & 0 \\
\end{pmatrix}$\\
Ну а звідси випливає, що:\\
$\mathbb{A} = (\mathbb{A} - \lambda \mathbb{I}) + \lambda \mathbb{I}$\\
Тут $\lambda \mathbb{I} = \begin{pmatrix}
\lambda & \dots & 0\\
\vdots & \ddots & \vdots\\
0 & \dots & \lambda
\end{pmatrix}$\\
$\Rightarrow \mathbb{A} = \begin{pmatrix}
\lambda & 1 & 0 & 0 & \dots & 0 \\
0 & \lambda & 1 & 0 & \dots & 0 \\
0 & 0 & \lambda & 1 & \dots & 0 \\
\vdots & \vdots & \vdots & \vdots & \ddots & \vdots \\
0 & 0 & 0 & 0 & \dots & 1 \\
0 & 0 & 0 & 0 & \dots & \lambda \\
\end{pmatrix} \overset{\textrm{позн.}}{=} J(\lambda)$\\
Таку матрицю називають \textbf{клітиною Жордана}
\fi

Розглянемо особливий випадок, коли $A: L \to L$ має єдиний (взагалі єдиний л.н.з.) власний вектор $f$, що відповідає власному числу $\lambda$. При цьому $\dim L = n$. Це один з випадків, де оператор $A$ не може бути діагоналізованим.\\
Наведемо декілька лем:

\begin{lemma}
$\ker(A - \lambda I), \Im(A-\lambda I)$ - інваріантні підпростори відносно оператора $A$.
\end{lemma}

\begin{proof}
За вимогою, $\ker(A-\lambda I) = L_{\lambda} = \linspan\{f\}$.\\
Тоді $\forall g \in \linspan\{f\}: g = \alpha f \implies Ag = A(\alpha f) = \alpha Af = (\alpha \lambda) f \implies Ag \in \linspan\{f\}$. Або теж саме, що $\forall g \in (\ker(A - \lambda I)) \implies Ag \in  \ker(A - \lambda I))$.
\bigskip \\
Тоді $\forall y \in \Im(A-\lambda I): \exists x \in L: y = (A-\lambda I)x \implies \\ Ay = A(A-\lambda I)x = (A^2-\lambda A)x = (A-\lambda I)(Ax) \in \Im(A-\lambda I)$.\\
Отже, $\ker(A-\lambda I)$ та $\Im(A-\lambda I)$ - два інваріантних підпростори.
\end{proof}

\begin{definition}
Задано $A: L \to L$ - лінійний оператор, в якому $f$ - власний вектор з власним числом $\lambda$.\\
Елемент $h \in L$ називається \textbf{приєднаним вектором} до власного вектора $f$ (висоти 1), якщо
\begin{align*}
(A-\lambda I)h = f
\end{align*}
\end{definition}

\begin{lemma}
Для $f$ існує приєднаний власний вектор $h$.
\end{lemma}

\begin{proof}
За вимогою, $\dim (\ker(A - \lambda I)) = \dim L_{\lambda} = 1$. Отже, за рівністю про зв'язок з ядром та образом, $\dim (\Im (A - \lambda I)) = n-1$.\\
Задамо таку множину $M=\ker(A-\lambda I) \cap \Im(A-\lambda I)$. Оскільки $M \subset \ker (A-\lambda I)$, то за одною лемою, $\dim M \leq \dim(\ker (A-\lambda I)) = 1$.
Тоді звідси або $\dim M = 0$, або $\dim M = 1$.
\bigskip \\
Розглянемо випадок $\dim M = 0$. Звідси маємо $M = \ker(A-\lambda I) \cap \Im(A-\lambda I) = \{0\}$.\\
Розглянемо $A |_{\ker(A-\lambda I)}: \ker(A-\lambda I) \to \ker(A-\lambda I)$ - звужений оператор. У нього є власний вектор $f$, оскільки $f \in \ker (A-\lambda I)$.\\
Також розглянемо $A |_{\Im(A-\lambda I)}: \Im(A-\lambda I) \to \Im(A-\lambda I)$ - звужений оператор. У цього оператора є власне число $\mu$ та власний вектор $g$ (згідно з першим зауваженням цього підрозділу):\\
$Ag = A |_{\Im(A-\lambda I)} g = \mu g$\\
Таким чином, $f \in \ker(A-\lambda I)$ та $g \in \Im(A-\lambda I)$. Оскільки $\ker ( A-\lambda I ) \cap \Im ( A - \lambda I ) = \{ 0 \}$, то вони утворюють пряму суму, а водночас $\{f,g\}$ - л.н.з. і є власними векторами для $A$ - суперечність. Отже, $\dim M \neq 0$.
\bigskip \\
Залишається єдиний випадок - це $\dim M = 1$.\\
$\begin{cases}
M \subset \ker(A-\lambda I)\\
\dim{(\ker(A-\lambda I))} = 1
\end{cases}
\implies M = \ker(A-\lambda I)$.
\bigskip \\
Тоді як можна побачити, $\ker(A-\lambda I) = \ker(A-\lambda I) \cap \Im (A - \lambda I)$, тобто $\ker(A-\lambda I) \subset \Im(A-\lambda I)$.\\
$f \in \ker (A- \lambda I) \implies f \in \Im(A - \lambda I)$.\\
Тоді $\exists h \in L: f = (A-\lambda I)h$, тобто знайшли приєднаний вектор.
\end{proof}

\begin{corollary}
Виділю окремо отриманий результат: $\ker (A-\lambda I) \subset \Im (A - \lambda I)$.
\end{corollary}

\begin{definition}
Задано $A: L \to L$ - лінійний оператор, в якому $f$ - власний вектор з власним числом $\lambda$.\\
Елемент $h^{(k)} \in L$ називається \textbf{приєднаним вектором} до власного вектора $f$ \textbf{висоти $k$}, якщо
\begin{align*}
(A-\lambda I)h^{(k)} = h^{(k-1)}
\end{align*}
\end{definition}

\begin{remark}
Буду позначати $h^{(1)} = h$.
\end{remark}

\begin{remark}
Рівняння для приєднаного висоти $k$ можна записати таким чином:\\
$(A-\lambda I)f = 0$\\
$(A-\lambda I)h^{(1)} = f \implies (A-\lambda I)^2 h^{(1)} = 0$\\
$(A-\lambda I)h^{(2)} = h \implies (A-\lambda I)^3 h^{(2)} = 0$\\
$\vdots$\\
$(A-\lambda I)^{k+1}h^{(k)} = 0$\\
Остаточно отримаємо таку форму:
\begin{align*}
(A-\lambda I)^{k+1}h^{(k)} = 0
\end{align*}
Взагалі-то кажучи, можна це проілюструвати ось так:
\begin{figure}[H]
\centering
\begin{tikzcd}
0 & f \arrow{l}{A-\lambda I} & h^{(1)} \arrow{l}{A-\lambda I} & h^{(2)} \arrow{l}{A-\lambda I} & \dots \arrow{l}{A-\lambda I} & h^{(k)} \arrow{l}{A-\lambda I}
\end{tikzcd}
\end{figure}
\end{remark}

\begin{lemma}
Нехай $\{f,h,h^{(2)},\dots,h^{(k)}\}$ - ланцюг власного та приєднаних до нього векторів.\\
Тоді $\{f,h,h^{(2)},\dots,h^{(k)}\}$ - л.н.з.
\end{lemma}

\begin{proof}
$\alpha_0 f + \alpha_1 h^{(1)} + \dots + \alpha_k h^{(k)} = 0$.\\
Подіємо оператором $(A-\lambda I)$ покроково $k$ разів:\\
$(A-\lambda I)(\alpha_0 f + \alpha_1 h^{(1)} + \dots + \alpha_k h^{(k)}) = 0$\\
$\alpha_1 f + \alpha_2 h^{(1)} + \dots + \alpha_k h^{(k-1)} = 0$\\
$\alpha_2 f + \dots + \alpha_k h^{(k-2)} = 0$\\
$\vdots$\\
$\alpha_{k-1} f + \alpha_k h^{(1)} = 0$\\
$\alpha_k f = 0$\\
$\Rightarrow \alpha_k = 0 \Rightarrow \alpha_{k-1} = 0 \Rightarrow \dots \Rightarrow \alpha_2 = 0 \Rightarrow \alpha_1 = 0 \Rightarrow \alpha_0 = 0$.\\
Отже, система $\{f,h,h^{(2)},\dots,h^{(k)}\}$ - л.н.з.
\end{proof}

\begin{theorem}
Задано $A: L \to L$ - лінійний оператор, в якому $f$ - єдиний л.н.з. власний вектор, $\dim L = n$. Тоді для власного вектора $f$ знайдеться ланцюг власного та приєднаних до нього векторів $\{f,h,h^{(2)},\dots,h^{(n-1)}\}$.
\end{theorem}

\begin{proof}
Спочатку доведемо лему:
\begin{lemma}
Позначимо $L_k = \Im (A-\lambda I)^k$. Тоді $L \supset L_1 \supset L_2 \supset \dots \supset L_{n-1} \supset \{0\}$.
\iffalse \textit{Дана лема притаманна будь-якому оператору $A: L \to L$.} \fi
\end{lemma}

\begin{proof}
Розглянемо оператор $A_1 = A |_{L_1}: L_1 \to L_1$ ($L_1$ - інваріантний за лемою).\\
У нього єдиний л.н.з. власний вектор $f$, оскільки $\ker(A-\lambda I) \subset L_1 \subset L$. \\
Тоді $\ker(A_1-\lambda I) = \ker(A-\lambda I) = \linspan\{f\}$. \\ Тому $\dim(\ker(A_1 - \lambda I)) = 1 \implies \dim (\Im(A_1-\lambda I)) = (n-1)-1=n-2$.\\
$\Im(A_1-\lambda I)$ - інваріантний підпростір для $A_1$, а отже, й для $A$ за лемою.\\
Зауважимо, що $\forall y \in L_2: \exists z \in L_1: y = (A_1 - \lambda I)z = (A-\lambda I)z \boxed{=}$\\
$z \in L_1 = \Im(A-\lambda I) \Rightarrow \exists x \in L: z = (A-\lambda I)x$\\
$\boxed{=} (A-\lambda I)^2 x$.\\
Отримали, що $\Im(A_1 - \lambda I) = \Im(A-\lambda I)^2 = L_2$, $\dim L_2 = n-2$.
\bigskip \\
Розглянемо оператор $A_2 = A |_{L_2}: L_2 \to L_2$.\\
У нього єдиний л.н.з. власний вектор $f$ за аналогічними міркуваннями.\\
Тоді $\ker(A_2-\lambda I) = \ker(A-\lambda I) = \linspan\{f\}$ \\
Тому $\dim(\ker(A_2 - \lambda I)) = 1 \implies \dim (\Im(A_2-\lambda I)) = (n-2)-1=n-3$.\\
$\Im(A_2-\lambda I)$ - інваріантний підпростір для $A_2$, а отже, й для $A$. Отримаємо, що $\dim (A_2 - \lambda I) = \dim (A-\lambda I)^3$. І знову теж саме...
\bigskip \\
Тоді остаточно, $L \supset \underset{=\Im(A-\lambda I)}{L_1} \supset \underset{=\Im(A-\lambda I)^2}{L_2} \supset \dots \supset \underset{=\Im(A-\lambda I)^{n-1}}{L_{n-1}} \supset \{0\}$.
\end{proof}
Додатково зауважимо, що
$\begin{cases}
\dim (\Im(A-\lambda I)^{n-1}) = 1\\
\dim (\ker (A-\lambda I)) = 1\\
\ker (A-\lambda I) \subset \Im(A-\lambda I)^{n-1}
\end{cases} \implies \Im(A-\lambda I)^{n-1} = \ker(A-\lambda I)
$.\\
Знайдемо ланцюг власного та приєднаного векторів довжини $n$.\\
Маємо $f$ - власний вектор.\\
$f \in L_{n-1} \implies f \in L_1 \Rightarrow \exists h^{(1)} \in L_1 \Rightarrow h^{(1)} \in L: (A-\lambda I)h = f$\\
$f \in L_{n-1} \implies f \in L_2 \Rightarrow \exists h^{(2)} \in L_2 \Rightarrow h^{(2)} \in L: (A-\lambda I)^2 h^{(2)} = f$\\
$\vdots$\\
$f \in L_{n-1} \implies \exists h^{(n-1)} \in L_{n-1} \Rightarrow h^{(n-1)} \in L: (A-\lambda I)^{(n-1)}h^{(n-1)} = f$.\\
Таким чином, ми знайшли ланцюг \begin{tikzcd}
0 & f \arrow{l}{A-\lambda I} & h^{(1)} \arrow{l}{A-\lambda I} & h^{(2)} \arrow{l}{A-\lambda I} & \dots \arrow{l}{A-\lambda I} & h^{(n-1)} \arrow{l}{A-\lambda I}
\end{tikzcd}.
\end{proof}

\begin{theorem}
Задано $A: L \to L$ - лінійний оператор, в якому $f$ - єдиний власний вектор з власним числом $\lambda$ та нехай $\dim L = n$. Відомо, що $\{f,h,h^{(2)},\dots,h^{(n-1)}\}$ - ланцюг власного та приєднаних до нього векторів.
Тоді $\{f,h,\dots,h^{(n-1)}\}$ - базис в просторі $L$.\\
\textit{Зрозуміло.}
\end{theorem}

Тоді можемо отримати матрицю оператора $A: L \to L$ в базисі $\{f,h,h^{(2)},\dots,h^{(n-1)}\}$.\\
$(A-\lambda I)f = 0$\\
$(A-\lambda I)h = f = f + 0h + \dots + 0h^{(k)}$\\
$(A-\lambda I)h^{(2)} = h = 0f + h + \dots + 0h^{(k)}$\\
$\vdots$\\
$(A-\lambda I)h^{(n-1)} = h^{(n-2)} = 0f + 0h + \dots + h^{(n-2)} + 0h^{(n-1)}$\\
$(\mathbb{A}- \lambda \mathbb{I}) = \begin{pmatrix}
0 & 1 & 0 & 0 & \dots & 0 \\
0 & 0 & 1 & 0 & \dots & 0 \\
0 & 0 & 0 & 1 & \dots & 0 \\
\vdots & \vdots & \vdots & \vdots & \ddots & \vdots \\
0 & 0 & 0 & 0 & \dots & 1 \\
0 & 0 & 0 & 0 & \dots & 0 \\
\end{pmatrix}$\\
Ну а звідси випливає, що $\mathbb{A} = (\mathbb{A} - \lambda \mathbb{I}) + \lambda \mathbb{I}$. Тут $\lambda \mathbb{I} = \begin{pmatrix}
\lambda & \dots & 0\\
\vdots & \ddots & \vdots\\
0 & \dots & \lambda
\end{pmatrix}$.\\
$\implies \mathbb{A} = \begin{pmatrix}
\lambda & 1 & 0 & 0 & \dots & 0 \\
0 & \lambda & 1 & 0 & \dots & 0 \\
0 & 0 & \lambda & 1 & \dots & 0 \\
\vdots & \vdots & \vdots & \vdots & \ddots & \vdots \\
0 & 0 & 0 & 0 & \dots & 1 \\
0 & 0 & 0 & 0 & \dots & \lambda \\
\end{pmatrix} \overset{\textrm{позн.}}{=} J(\lambda)$ - так звана \textbf{клітина Жордана}.

\subsection{Теорема Жордана}
\begin{theorem}[Теорема Жордана]
Задано $A: L \to L$ - лінійний оператор. Тоді в $L$ є базис з власних та приєднаних векторів:\\
$\{f_1, h_1^1, \dots ,h_1^{k_1}, \hspace{1cm} f_2, h_2^1, \dots, h_2^{k_2}, \hspace{1cm} \dots \hspace{1cm} f_m, h_m^{1}, \dots, h_m^{k_m}\}$.
\end{theorem}

\begin{proof}
Доведення за індукцією за $\dim L$. Але перед цим маємо оператор $A: L \to L$ - у нього існує деяке власне число $\lambda$. Тому в подальшому розглядатимемо оператор $B = A -\lambda I: L \to L$.\\
База індукції: $\dim L = 1$.\\
Тоді $L = \linspan\{f\}$, $Af \in L \implies Af = \alpha f \implies f$ - власний вектор.
\bigskip \\
Ще одна база індуції: $\dim L = 2$.\\
Тоді $\exists f_1: Af = \lambda f_1$. А ось далі є два варіанти:\\
1) $\exists f_2$ - другий власний вектор, щоб $\{f_1,f_2\}$ були л.н.з. Тоді $\{f_1,f_2\}$ - базис одразу.\\
2) $f_1$ - єдиний л.н.з. власний вектор, тоді за попередньою теоремою, існує приєднаний вектор $h_1$. Тоді $\{f,h_1\}$ - базис.\\
\\
Крок індукції: нехай твердження виконується для $\dim L < n$.\\
Доведемо існування базису для випадку $\dim L = n$.\\
Оскільки для оператора $B$ існує власне число $\lambda$, то звідси $\ker B \neq \{0\}$, а тому \\ $\dim \Im B = \dim L - \dim \ker B < n$.\\
\iffalse
Зауважимо, що $\Im B$ - інваріантний простір відносно оператора $A$. Справді,\\
$\forall y \in \Im B \implies y = Bx \implies Ay = ABx = A(A-\lambda I)x = (A^2-\lambda A)x = (A-\lambda I) (Ax) = B(Ax) \in \Im B$.\\
А тому зараз розглянемо звужений оператор $A |_{\Im B}: \Im B\to \Im B$.\\
\fi
Зауважимо, що $\Im B$ - інваріантний відносно оператора $B$ (уже доводили). А тому розглянемо звужений оператор $B |_{\Im B}: \Im B\to \Im B$.\\
За припущенням індукції, в $\Im B$ існує базис з власних та приєднаних, оскільки $\dim \Im B < n$.\\
$ \lambda: \begin{cases}
\begin{tikzcd}
0 & f_1 \arrow{l}{B} & h_1^1 \arrow{l}{B} & \dots \arrow{l}{B} & h_1^{k_1} \arrow{l}{B}
\end{tikzcd}\\
\vdots \\
\begin{tikzcd}
0 & f_s \arrow{l}{B} & h_s^1 \arrow{l}{B} & \dots \arrow{l}{B} & h_s^{k_s} \arrow{l}{B}
\end{tikzcd}
\end{cases}
\\
\lambda_{s+1}: f_{s+1}, h_{s+1}^1,\dots,h_{s+1}^{k_{s+1}}\\
\vdots\\
\lambda_m: f_{m}, h_{m}^1,\dots,h_{m}^{k_{m}}\\
$
Це все - базис в $\Im B$. Важливо, що в $\lambda_{s+1},\dots,\lambda_m$ власні числа можуть співпадати, але жодна з цих не дорівнює $\lambda$. Також зауважу, що може навіть таке бути, що нема ланцюга, що відповідає $\lambda$, або абсолютно всі башти відповідають власному числу $\lambda$.\\
Тепер нам необхідно розширити до базису $L$, але яким чином це відбувається, зараз дізнаємось. Заздалегідь позначу $L_j = \linspan\{f_j,h_j^{1},\dots,h_j^{k_j}\}$, де $j=\overline{1,s} \cup \overline{s+1,m}$.

\begin{lemma}
$L_j$ - інваріантні підпростори відносно $B$.
\end{lemma}

\begin{proof}
Позначу $x = \alpha_j^0 f_j + \alpha_j^1 h_j^1 + \dots + \alpha_j^{k_j} h_j^{k_j}$.\\
При $j=\overline{1,s}$ якщо $x \in L_j$, то звідси\\
$Bx = (A-\lambda I)x = \alpha_j^1 f_j + \alpha_j^2 h_j^1 + \dots + \alpha_j^{k_j}h_j^{k_j-1} \implies Bx \in L_j$.
\bigskip \\
При $j=\overline{s+1,m}$ маємо\\
$Bf_j = (A-\lambda I)f_j = Af_j - \lambda f_j = (\lambda_j - \lambda)f_j$;\\
$Bh_j^1 = (A-\lambda I)h_j^1 = (A-\lambda_j I+\lambda_j I-\lambda I)h_j^1 = (A-\lambda_jI)h_j^1 + (\lambda_jI-\lambda I)h_j^1 = f_j + (\lambda_j-\lambda)h_j^1$;\\
$\vdots$\\
$Bh_j^{k_j} = (A-\lambda I)h_j^{k_j} = (A-\lambda_j I+\lambda_j I-\lambda I)h_j^{k_j} = (A-\lambda_jI)h_j^{k_j} + (\lambda_jI-\lambda I)h_j^{k_j} = h_j^{k_j-1} + (\lambda_j-\lambda)h_j^{k_j}$.\\
Отримані $Bf_j,Bh_j^1,\dots,B_j^{k_j} \in L_j$, а тому отримаємо:\\
$Bx = \alpha_j^0 Bf_j + \alpha_j^1 Bh_j^1 + \dots + \alpha_j^{k_j}Bh_j^{k_j} \implies Bx \in L_j$.\\
Все, ми довели інваріантність $L_j$ відносно $B$.
\end{proof}

\iffalse
\begin{lemma}
$\Im B = L_1  \oplus  L_2  \oplus  \dots  \oplus  L_m$.
\end{lemma}

\begin{proof}
Покажемо, що $L_j \cap L_i = \{0\}, j \neq i$. Справді,\\
$z \in L_j \cap L_i \Rightarrow z = \left[\begin{gathered} z^0_j f_j + z^1_j h^1_j + \dots + z_j^{k_j} h_j^{k_j} \\ z_i^0f_i + z_i^1 h_i^1 + \dots + z_i^{k_i} h_i^{k_i} \end{gathered} \right.$\\
$\implies z^0_j f_j + z^1_j h^1_j + \dots + z_j^{k_j} h_j^{k_j} + (-z_i^0)f_i + (-z_i^1) h_i^1 + \dots + (-z_i^{k_i}) h_i^{k_i} = 0$\\
За побудовою, $\{L_j, L_i\}$ - елементи з базису $\Im B$, тому є л.н.з.\\
$\implies z^0_j = z^1_j = \dots = z^{k_j}_j = z^0_i = z^1_i = \dots = z^{k_i}_i = 0$\\
Отже, $z = 0$, а тому $L_j \cap L_i = \{0\}$.\\
Таким чином, отримали, що $\Im B = L_1  \oplus  L_2  \oplus  \dots  \oplus  L_m$.
\end{proof}
\fi

\begin{lemma}
Ланцюги, що відповідають власному числу $\lambda$, будуть подовженими на один новий приєднаний вектор із простору $L$. А ланцюги, що відповідають власному числу $\lambda_2,\dots,\lambda_m$, не будуть подовженими.
\end{lemma}

\begin{proof}
А тепер розглянемо звужений оператор $B|_{L_j}: L_j \to L_j$. Знову розіб'ємо на два випадки.\\
При $j=\overline{1,s}$ зауважимо, що оскільки $B|_{L_j}f_j = 0$, то звідси $\ker B|_{L_j} \neq \{0\}$. Отже, оператор $B|_{L_j}$ не має оберненого. Тоді елемент $h_j^{k_j}$ не є прообразом жодного елементу в просторі $L_j$ (!)\\
\textit{Інакше якщо припустити, що $h_j^{k_j}$ є прообразом деякого елемента $x \in L_j$, тобто щоб $B|_{L_j}h_j^{k_j} = x$, то ми прийдемо факту, що $\{f_j,h_j^1,\dots,h_j^{k_j}\}$ буде л.з. що суперечить попередньому підрозділу.}\\
Оскільки $h_j^{k_j} \in \Im B$, просто тому що $L_j \subset \Im B$, то ми можемо знайти елемент $h_j^{k_j+1}$, що лежить за межами $L_j$, власне $h_j^{k_j+1} \in L$, щоб $Bh_j^{k_j+1} = h_j^{k_j}$. Тобто ми змогли подовжити ланцюг.\\
\textit{Чи можемо ми подовжити один ланцюг на два нових приєднаних або більше? Ні. Можна припустити, що $Bh_j^{k_j+2} = h_j^{k_j+1}$, тобто є друге подовження, але тоді $h_j^{k_j+1} \in \Im B$, ми швиденько прийдемо до суперечності.}
\bigskip \\
При $j=\overline{s+1,m}$ спробуємо довести, що $\ker B|_{L_j} = \{0\}$.
\iffalse
$x \in \ker B|{L_j} \implies B|_{L_j} x = 0 \implies Bx = 0 \implies Ax = \lambda x$, тобто $x$ - власний вектор для числа $\lambda$. А оскільки $L_j$ не містить власних векторів для $\lambda$, то звідси $x = 0$.\\
Це означає, що існує $(B|_{L_j})^{-1}$. А тому будь-який елемент $L_j$ при дії $B|_{L_j}$ має прообраз з $L_j$. А це означає, що $h_{s+1}^{k_{s+1}},\dots,h_m^{k_m}$ не мають продовження за межами $L_{s+1},\dots,L_m$.
\fi
Якщо $x \in \ker B|_{L_j}$, то звідси $B|_{L_j}x = 0$.\\
Але оскільки $\ker B|_{L_j} \subset L_j$, то звідси $x \in L_j \implies x = \alpha_j^0 f_j + \alpha_j^1 h_j^1 + \dots + \alpha_j^{k_j}h_j^{k_j}$.\\
$B|_{L_j}x = \alpha_j^0 B|_{L_j}f_j + \alpha_j^1 B|_{L_j}h_j^1 + \dots + \alpha_j^{k_j} B|_{L_j}h_j^{k_j} = 0$ (*).\\
Нагадаю з попередньої леми, що в цьому випадку $B|_{L_j}f_j = (\lambda_j-\lambda)f_j$, $B|_{L_j}h_j^1 = f_j + (\lambda_j - \lambda)h_j^1$, \dots $B|_{L_j}h_j^{k_j} = h_j^{k_j-1} + (\lambda_j - \lambda)h_j^{k_j}$. Я позначу $\lambda-\lambda_j = \mu_j \neq 0$. Підставимо це все в наш вираз (*):
$\alpha_j^0 \mu_j f_j + \alpha_j^1 (f_j+\mu_j h_j^1) + \dots + \alpha_j^{k_j} (h_j^{k_j-1}+\mu_j h_j^{k_j}) = \\ =(\alpha_j^0 \mu_j + \alpha_j^1) f_j + (\alpha_j^1 \mu_j + \alpha_j^2) h_j^1 + \dots + (\alpha_j^{k_j-1}\mu_j + \alpha_j^{k_j})h_j^{k_j-1} + \alpha_j^{k_j} \mu_j h_j^{k_j} = 0$.\\
У силу л.н.з. системи з $L_j$ ми маємо $\alpha_j^{k_j} = 0 \implies \alpha_j^{k_j-1} = 0 \implies \dots \implies \alpha_j^0 = 0$.\\
Підставимо ці коефіцієнти в $x$ та отримаємо $x = 0$. Отже, $\ker B|_{L_j} = \{0\}$.\\
Це означає, що існує $(B|_{L_j})^{-1}$. А тому будь-який елемент $L_j$ при дії $B|_{L_j}$ має прообраз з $L_j$. А це означає, що $h_{s+1}^{k_{s+1}},\dots,h_m^{k_m}$ не мають продовження за межами $L_{j}$.
\end{proof}

Додатково до цього ми маємо таку штуку. Оскільки $\ker B|_{\Im B} = \linspan \{f_1,\dots,f_s\}$ та $\ker B \supset \ker B|_{\Im B}$, то система $\{f_1,\dots,f_s\}$ в просторі $\ker B$ є л.н.з., тому розширимо до базису $\{f_1,\dots,f_s,g_1,\dots,g_r\}$.\\
Навіщо нам ці $g_1,\dots,g_r$? Це ті власні вектори, що відповідають власному числу $\lambda$, але вони не потрапили до $\Im B$. Бо в базисі $\Im B$, нагадаю, може й не бути ланцюга з власним числом $\lambda$.
\bigskip \\
Таким чином, ми розширили систему в лінійному просторі $L$ до такой системи:\\
$ \lambda: \begin{cases}
\begin{tikzcd}
0 & f_1 \arrow{l}{B} & h_1^1 \arrow{l}{B} & \dots \arrow{l}{B} & h_1^{k_1} \arrow{l}{B} & \textcolor{red}{h_1^{k_1+1}} \arrow{l}{B}
\end{tikzcd}\\
\vdots \\
\begin{tikzcd}
0 & f_s \arrow{l}{B} & h_s^1 \arrow{l}{B} & \dots \arrow{l}{B} & h_s^{k_s} \arrow{l}{B} & \textcolor{red}{h_s^{k_s+1}} \arrow{l}{B}
\end{tikzcd}\\
\begin{tikzcd}
0 & \textcolor{red}{g_1} \arrow{l}{B} & \dots & 0 & \textcolor{red}{g_r} \arrow{l}{B} &
\end{tikzcd}
\end{cases}
\\
\lambda_{s+1}: f_{s+1}, h_{s+1}^1,\dots,h_{s+1}^{k_{s+1}}\\
\vdots\\
\lambda_m: f_{m}, h_{m}^1,\dots,h_{m}^{k_{m}}
$.\\
Нарешті, лишилось довести таку лему:
\begin{lemma}
Система $\{f_1,h_1^1,\dots,h_1^{k_1},\textcolor{red}{h_1^{k_1+1}}, \hspace{0.5cm} \dots \hspace{0.5cm} f_s, h_s^{1}, \dots,h_s^{k_s}, \textcolor{red}{h_s^{k_s+1}}, \\ f_{s+1}, h_{s+1}^1, \dots, h_{s+1}^{k_{s+1}} \hspace{0.5cm} \dots \hspace{0.5cm} f_{m}, h_{m}^1, \dots, h_{m}^{k_m}, \hspace{0.5cm} \textcolor{red}{g_1, \dots, g_r}\}$ утворює базис в лінійному просторі $L$.
\end{lemma}

\begin{proof}
I. Перевіримо на л.н.з.\\
$\left(\alpha_1^0 f_1 + \alpha_1^1 h_1^1 + \dots + \alpha_1^{k_1}h_1^{k_1} + \alpha_1^{k_1+1} h_1^{k_1+1}\right) + \dots + \left(\alpha_s^0 f_s + \alpha_s^1 h_s^1 + \dots + \alpha_s^{k_s}h_s^{k_s} + \alpha_s^{k_s+1} h_s^{k_s+1} \right) + \\
+\left(\alpha_{s+1}^0 f_{s+1} + \alpha_{s+1}^1 h_{s+1}^1 + \dots + \alpha_{s+1}^{k_{s+1}}h_{s+1}^{k_{s+1}} \right) + \dots + \left(\alpha_m^0 f_m + \alpha_m^1 h_m^1 + \dots + \alpha_m^{k_m}h_m^{k_m}\right) + \\ +\beta_1 g_1 + \dots + \beta_r g_r = 0$\\
Подіємо це чудо оператором $B = A - \lambda I$ - отримаємо:\\
$\left(\alpha_1^0 0 + \alpha_1^1 f_1 + \dots + \alpha_1^{k_1}h_1^{k_1-1} + \alpha_1^{k_1+1} h_1^{k_1}\right) + \dots + \left(\alpha_s^0 0 + \alpha_s^1 f_1 + \dots + \alpha_s^{k_s}h_s^{k_s-1} + \alpha_s^{k_s+1} h_s^{k_s}\right) + \\
B\left(\alpha_{s+1}^0 f_{s+1} + \alpha_{s+1}^1 h_{s+1}^1 + \dots + \alpha_{s+1}^{k_{s+1}}h_{s+1}^{k_{s+1}} \right) + \dots + B\left(\alpha_m^0 f_m + \alpha_m^1 h_m^1 + \dots + \alpha_m^{k_m}h_m^{k_m}\right) + \\
+B(\beta_1 g_1 + \dots + \beta_r g_r) = 0$\\
У нас $B\left(\alpha_{s+1}^0 f_{s+1} + \alpha_{s+1}^1 h_{s+1}^1 + \dots + \alpha_{s+1}^{k_{s+1}} h_{s+1}^{k_{s+1}} \right) \in L_{s+1}$, звідси можна розкласти за базисом в $L_{s+1}$, тобто\\
$B\left(\alpha_{s+1}^0 f_{s+1} + \alpha_{s+1}^1 h_{s+1}^1 + \dots + \alpha_{s+1}^{k_{s+1}} h_{s+1}^{k_{s+1}} \right) = \widetilde{\alpha_{s+1}^0}f_{s+1} + \widetilde{\alpha_{s+1}^1}f_{s+1}^1+\dots + \widetilde{\alpha_{s+1}^{k_{s+1}}}h_{s+1}^{k_{s+1}}$;\\
\vdots \\
У нас $B\left(\alpha_{m}^0 f_{m} + \alpha_{m}^1 h_{m}^1 + \dots + \alpha_{m}^{k_m} h_{m}^{k_m} \right) \in L_{m}$, звідси можна розкласти за базисом в $L_{m}$, тобто\\
$B\left(\alpha_{m}^0 f_{m} + \alpha_{m}^1 h_{m}^1 + \dots + \alpha_{m}^{k_m} h_{m}^{k_m} \right) = \widetilde{\alpha_{m}^0}f_{m} + \widetilde{\alpha_{m}^1}f_{m}^1+\dots + \widetilde{\alpha_{m}^{k+1}}h_{m}^{k_{m}}$.\\
Нарешті, $B(\beta_1 g_1 + \dots + \beta_r g_r) = 0$, оскільки $g_1,\dots,g_r \in \ker B$.\\
Отже, матимемо таке рівняння:\\
$\left(\alpha_1^0 0 + \alpha_1^1 f_1 + \dots + \alpha_1^{k_1}h_1^{k_1-1} + \alpha_1^{k_1+1} h_1^{k_1}\right) + \dots + \left(\alpha_s^0 0 + \alpha_s^1 f_1 + \dots + \alpha_s^{k_s}h_s^{k_s-1} + \alpha_s^{k_s+1} h_s^{k_s}\right) + \\
\left(\widetilde{\alpha_{s+1}^0}f_{s+1} + \widetilde{\alpha_{s+1}^1}f_{s+1}^1+\dots + \widetilde{\alpha_{s+1}^{k_{s+1}}}h_{s+1}^{k_{s+1}}\right) + \dots + \left(\widetilde{\alpha_{m}^0}f_{m} + \widetilde{\alpha_{m}^1}f_{m}^1+\dots + \widetilde{\alpha_{m}^{k+1}}h_{m}^{k_{m}}\right) + 0 = 0$.\\
Оскільки $\{f_1, h_1^1, \dots ,h_1^{k_1}, \hspace{0.5cm} f_2, h_2^1, \dots, h_2^{k_2}, \hspace{0.5cm} \dots \hspace{0.5cm} f_m, h_m^{1}, \dots, h_m^{k_m}\}$ - базиси $\Im B$, то вона л.н.з. в просторі $L$, звідси\\
$\alpha_1^1 = \dots = \alpha_1^{k_1} = \alpha_1^{k_1+1} = \dots = \alpha_s^1 = \dots = \alpha_s^{k_s} = \alpha_s^{k_s+1} = \\
=\widetilde{\alpha_{s+1}^0} = \widetilde{\alpha_{s+1}^1} = \dots = \widetilde{\alpha_{s+1}^{k_{s+1}}} = \dots = \widetilde{\alpha_m^0} = \widetilde{\alpha_m^1} = \dots = \widetilde{\alpha_m^{k_m}} = 0$\\
Підставимо ці коефіцієнти в початкове рівняння - отримаємо:\\
$\alpha_1^0f_1 + \dots + \alpha_s^0 f_s + \beta_1 g_1 + \dots + \beta_r g_r = 0$.\\
А оскільки $\{f_1,\dots,f_s,g_1,\dots,g_r\}$ - базис $\ker B$, тоді вона л.н.з. в просторі $L$, тоді миттєво \\ $\alpha_1^0 = \dots = \alpha_s^0 = \beta_1 = \dots = \beta_r = 0$.\\
Отже, наша початкова система - л.н.з.
\bigskip \\
II. Перевіримо на повноту.\\
$\forall z \in L: Bz \in \Im B \implies \\ Bz = \alpha_1^0 f_1 + \alpha_1^1 h_1^1 + \dots + \alpha_1^{k_1}h_1^{k_1} + \dots + \alpha_s^0 f_s + \alpha_s^1 h_s^1 + \dots + \alpha_s^{k_s} h_s^{k_s} + \\ + \alpha_{s+1}^0 f_{s+1} + \alpha_{s+1}^1 h_{s+1}^1 + \dots + \alpha_{s+1}^{k_{s+1}}h_{s+1}^{k_{s+1}} + \dots + \alpha_m^0 f_m + \alpha_m^1 h_m^1 + \dots + \alpha_m^{k_m}h_m^{k_m}$.\\
Розглянемо елемент $w \in L$, що має таку формулу:\\
$w = \alpha_1^0 h_1^1 + \alpha_1^1 h_1^2 + \dots + \alpha_1^{k_1}h_1^{k_1+1} + \dots + \alpha_s^0 h_s^1 + \alpha_s^1 h_s^2 + \dots + \alpha_s^{k_s}h_s^{k_s+1} + \\
+ \widetilde{\alpha_{s+1}^0} f_{s+1} + \widetilde{\alpha_{s+1}^1} h_{s+1}^1 + \dots + \widetilde{\alpha_{s+1}^{k_{s+1}}} h_{s+1}^{k_{s+1}} + \dots + \widetilde{\alpha_{m}^0} f_{m} + \widetilde{\alpha_{m}^1} h_{m}^1 + \dots + \widetilde{\alpha_{m}^{k_{m}}} h_{m}^{k_{m}}$.\\
Спеціально ми так підібрали, щоб подіявши оператором $B$, ми могли отримати елемент $z$.\\
Тут $\widetilde{\alpha}$ підібрані так, що:\\
$B(\widetilde{\alpha_j^0} f_j + \dots + \widetilde{\alpha_j^{k_j}} h_j^{k_j}) = \alpha_j^0 f_j + \dots + \alpha_j^{k_j}h_j^{k_j}, j = \{s+1, \dots, m\}$.\\
$\implies B(z-w) = Bz - Bw = 0$, отже, $z-w \in \ker B$.\\
$\implies z = w + \tau_1 f_1 + \dots + \tau_s f_s + \gamma_1 g_1 + \dots + \gamma_r g_r$.\\
Завершальний етап - підставити $w$ в наш вираз:\\
$z = \left(\tau_1 f_1 + \alpha_1^0 h_1^1 + \alpha_1^1 h_1^2 + \dots + \alpha_1^{k_1}h_1^{k_1+1}\right) + \dots + \left(\tau_s f_s + \alpha_s^0 h_s^1 + \alpha_s^1 h_s^2 + \dots + \alpha_s^{k_s}h_s^{k_s+1}\right) + \\
+ \left(\widetilde{\alpha_{s+1}^0} f_{s+1} + \widetilde{\alpha_{s+1}^1} h_{s+1}^1 + \dots + \widetilde{\alpha_{s+1}^{k_{s+1}}} h_{s+1}^{k_{s+1}}\right) + \dots + \left(\widetilde{\alpha_{m}^0} f_{m} + \widetilde{\alpha_{m}^1} h_{m}^1 + \dots + \widetilde{\alpha_{m}^{k_{m}}} h_{m}^{k_{m}} \right)+\\
+\gamma_1 g_1 + \dots + \gamma_r g_r$.\\
Тобто будь-який елемент $z \in L$ розклався як лінійна комбінація нашої системи.\\
Остаточно: наша початкова система утворює базис в $L$.
\end{proof}

Після доведення останньої леми, ми можемо стверджувати, що для випадку $\dim L = n$ знайшовся базис власних та приєднаних векторів.\\
МІ доведено (нарешті).
\end{proof}

\begin{remark}
Доведення повноти досить неприємна штука. Можна цього кроку уникнути таким чином.\\
Із самого початку ми мали $\dim \Im B \overset{\text{позн}}{=} r < n$. За МІ, там є $r$ елементів, що описують власні та приєднані до них вектори. Зауважимо, що $\dim \Im B|_{L_j} = k_j$, якщо $j = \overline{1,s}$, а також $\dim \Im B|_{L_j} = k_j + 1$, якщо $j = \overline{s+1,m}$. Простори $\Im B_{L_1},\dots, \Im B_{L_m}$ утворюють пряму суму, а тому $\dim B|_{\Im B} = k_1 + \dots + k_s + (k_{s+1}+1) + \dots + (k_m+1) = (k_1 + 1) + \dots + (k_s + 1) + (k_{s+1}+1) + \dots + (k_m+1) - s = r - s$. Таким чином, $\dim \ker B|_{\Im B} = s$, та $\ker B|_{\Im B} \subset \ker B$ має базис $\{f_1,\dots,f_s\}$, тому розширяється ця система до $\{f_1,\dots,f_s, g_1,\dots,g_p\}$, де число $p = n-r-s$.\\
Потім доводиться подовження деяких ланцюгів на одну висоту.\\
Сумарна кількість: $(k_1+2)+ \dots + (k_s+2) + (k_{s+1}+1) + \dots + (k_m+1) + p = s+r + (n-r-s) = n$.\\
Оскільки $\dim L = n$, то нам достатньо довести л.н.з.
\end{remark}

\begin{theorem}
Задано $A: L \to L$ - лінійний оператор, маємо базис власних та приєднаних $\{\underset{\lambda_1}{f_1, h_1^1, \dots ,h_1^{k_1}}, \hspace{1cm} \underset{\lambda_2}{f_2, h_2^1, \dots, h_2^{k_2}}, \hspace{1cm} \dots \hspace{1cm} \underset{\lambda_m}{f_m, h_m^{1}, \dots, h_m^{k_m}}\}$.\\
Позначимо $L_j = \linspan\{f_j, h_j^1,\dots, h_j^{k_j}\}$. Тоді $L_j$ - інваріантні відносно $A$ та $L = L_1  \oplus  L_2  \oplus  \dots  \oplus  L_m$.
\end{theorem}

\begin{proof}
Нагадаю, що $L_j$ - інваріантний відносно оператору $B = A-\lambda I$, де $\lambda$ - деяке власне число. Але зрозуміло, що $L_j$ - інваріантний відносно $A$, оскільки\\
$\forall x \in L_j: Bx \in L_j \implies Ax = (A-\lambda I)x + \lambda Ix = Bx + \lambda x \in L_j$.\\
Покажемо, що $L_j \cap L_i = \{0\}, j \neq i$. Справді,\\
$z \in L_j \cap L_i \Rightarrow z = \left[\begin{gathered} z^0_j f_j + z^1_j h^1_j + \dots + z_j^{k_j} h_j^{k_j} \\ z_i^0f_i + z_i^1 h_i^1 + \dots + z_i^{k_i} h_i^{k_i} \end{gathered} \right.$\\
$\implies z^0_j f_j + z^1_j h^1_j + \dots + z_j^{k_j} h_j^{k_j} + (-z_i^0)f_i + (-z_i^1) h_i^1 + \dots + (-z_i^{k_i}) h_i^{k_i} = 0$\\
За побудовою, $\{L_j, L_i\}$ - елементи з базису, тому є л.н.з.\\
$\implies z^0_j = z^1_j = \dots = z^{k_j}_j = z^0_i = z^1_i = \dots = z^{k_i}_i = 0$\\
Отже, $z = 0$, а тому $L_j \cap L_i = \{0\}$.\\
Таким чином, отримали, що $L = L_1  \oplus  L_2  \oplus  \dots  \oplus  L_m$.
\end{proof}

А тепер знову задамо $A: L \to L$ - лінійний оператор. За теоремою Жордана, в просторі $L$ є базис з власних приєднаних до них векторів.\\
$\{\underset{\lambda_1}{f_1, h_1^1, \dots ,h_1^{k_1}}, \hspace{1cm} \underset{\lambda_2}{f_2, h_2^1, \dots, h_2^{k_2}}, \hspace{1cm} \dots \hspace{1cm} \underset{\lambda_m}{f_m, h_m^{1}, \dots, h_m^{k_m}}\}$\\
Щойно довели, що $L = L_1  \oplus  \dots +  \oplus  L_m$. Тому матриця оператора $A$ в цьому базисі має вигляд:\\
$ \mathbb{A}_{J} =
  \tikz[baseline=(M.west)]{%
    \node[matrix of math nodes,matrix anchor=west,left delimiter=(,right delimiter=),ampersand replacement=\&] (M) {%
      \mathbb{A}_{|_{L_1}} \& \mathbb{O} \& \dots \& \mathbb{O} \\
      \mathbb{O} \& \mathbb{A}_{|_{L_2}} \& \dots \& \mathbb{O} \\
      \vdots \& \vdots \& \ddots \& \vdots \\
      \mathbb{O} \& \mathbb{O} \& \dots \& \mathbb{A}_{|_{L_m}} \\
    };
    \node[draw,fit=(M-1-1)(M-1-1),inner sep=-1pt] {};
    \node[draw,fit=(M-2-2)(M-2-2),inner sep=-1pt] {};
    \node[draw,fit=(M-4-4)(M-4-4),inner sep=-1pt] {};
  }
$\\ 
Розглянемо матриці звужених операторів $A|_{L_j}$, $j = 1,\dots,m$ в силу інваріантності.\\
$L_j = \linspan\{f_j,h_j^1,\dots,h_j^{k_j}\}$ - базис з одного власного вектора та приєднаних до нього. Тоді матриця $A|_{L_j}$ в цьому базисі має вигляд: $\mathbb{A}|_{L_j} = J(\lambda_j)$ - клітина Жордана.\\
Остаточний вигляд матриці:\\
$ \mathbb{A}_{J} =
  \tikz[baseline=(M.west)]{%
    \node[matrix of math nodes,matrix anchor=west,left delimiter=(,right delimiter=),ampersand replacement=\&] (M) {%
      J(\lambda_1) \& \mathbb{O} \& \dots \& \mathbb{O} \\
      \mathbb{O} \& J(\lambda_2) \& \dots \& \mathbb{O} \\
      \vdots \& \vdots \& \ddots \& \vdots \\
      \mathbb{O} \& \mathbb{O} \& \dots \& J(\lambda_m) \\
    };
    \node[draw,fit=(M-1-1)(M-1-1),inner sep=-1pt] {};
    \node[draw,fit=(M-2-2)(M-2-2),inner sep=-1pt] {};
    \node[draw,fit=(M-4-4)(M-4-4),inner sep=-1pt] {};
  }
$ - \textbf{Жорданова форма матриці}.
\\
\bigskip \\
\fi

\iffalse
\begin{example}
Маємо матрицю оператора $\mathbb{A} = \begin{pmatrix}
2 & 1 & -1 & 8 & -3 & 1 & -1 & 1 & -1 & 1 \\
0 & 2 & 0 & 7 & 5 & 1 & -1 & 1 & -1 & 1 \\
0 & 0 & 2 & 7 & 5 & 1 & -1 & 1 & -1 & 1 \\
0 & 0 & 0 & 2 & 0 & 1 & -1 & 1 & -1 & 1 \\
0 & 0 & 0 & 0 & 2 & 1 & -1 & 1 & -1 & 1 \\
0 & 0 & 0 & 0 & 0 & 1 & 1 & 1 & -1 & 0 \\
0 & 0 & 0 & 0 & 0 & 0 & 1 & 0 & 0 & 1 \\
0 & 0 & 0 & 0 & 0 & 0 & 0 & 0 & 1 & 0 \\
0 & 0 & 0 & 0 & 0 & 0 & 0 & -1 & 2 & 1 \\
0 & 0 & 0 & 0 & 0 & 0 & 0 & 0 & 0 & 1
\end{pmatrix}$ - страшна матриця. Знайдемо її Жорданову форму.\\
Зауважимо, що ми можемо застосувати формулу $\det \begin{pmatrix}
 \mathbb{D} & \vline & \mathbb{C} \\
 \hline
 \mathbb{O} & \vline & \mathbb{B}
\end{pmatrix} = \det \mathbb{A} \det \mathbb{B}$ під час обшуку власних чисел. Тоді отримаємо (пропущу деякі прості кроки обчислень):\\
$\det (\mathbb{A} - \lambda \mathbb{I}) = (\lambda - 2)^5 (\lambda-1)^5 = 0$. Отже, маємо $\lambda = 2, \lambda = 1$.
\end{example}
\fi

%GB approach
\iffalse
\subsection{Квадратичні форми}
\begin{definition}
Задано $(E,(\cdot,\cdot))$ - евклідів простір та $A: E \to E$ - лінійний оператор.\\
\textbf{Квадратичною формою} на $E$ називають таке відображення $S: E \to \mathbb{R} (\mathbb{C})$, що:
\begin{align*}
\forall x \in E: S(x) = (x,Ax)
\end{align*}
\end{definition}

\begin{proposition}[Поляризаційна тотожність]
Задано $(E,(\cdot,\cdot))$ - евклідів простір та $S: E \to \mathbb{R}(\mathbb{C})$ - квадратична форма. Тоді за квадратичною формою можна відновити півторалінійну форму $\varphi(x,y)$ такою формулою:\\
$4 \varphi(x,y) = S(x+y) - S(x-y) + i(S(x+iy)-S(x-iy))$.
\iffalse
В дійсному випадку $2\varphi(x,y) - 2\varphi(y,x) = S(x+y)- S(x-y)$.
\fi
\end{proposition}

\begin{proof}
\iffalse
1) Випадок $\mathbb{C}$\\
$S(x+y) - S(x-y) + i(S(x+iy)-S(x-iy)) = \\ \varphi(x+y, x+y) - \varphi(x-y,x-y) + i(\varphi(x+iy,x+iy)-\varphi(x-iy,x-iy)) = \\
= \varphi(x,x) + \varphi(x,y) + \varphi(y,x) + \varphi(y,y) - \\
- (\varphi(x,x)-\varphi(x,y)-\varphi(y,x)+\varphi(y,y)) + \\
+ i(\varphi(x,x)-i\varphi(x,y)-i\varphi(y,x)-i i\varphi(y,y)) - \\
- i(\varphi(x,x)+i\varphi(x,y)-i\varphi(y,x)-ii\varphi(y,y)) = \\
= 2 \varphi(x,y) + 2 \varphi(y,x) + 2 \varphi(x,y) - 2 \varphi(y,x) = 4 \varphi(x,y)$
\bigskip \\
2) Випадок $\mathbb{R}$ - аналогічно
\fi
Маємо $S(x) = (Ax,x)$ для деякого оператора $A$. Тоді\\
$S(x+y) - S(x-y) + i(S(x+iy)-S(x-iy)) = \\ 
= (x+y,A(x+y)) - (x-y,A(x-y)) + i[(x+iy,A(x+iy)) - (x-iy,A(x-iy))] = \\
= (x+y, Ax+Ay) - (x-y,Ax-Ay) + i[(x+iy, Ax+iAy) - (x-iy,Ax-iAy)] = \\
= (x,Ax)+(x,Ay)+(y,Ax)+(y,Ay) - (x,Ax) + (x,Ay) + (y,Ax) - (y,Ay) + \\
+i[(x,Ax)+ \overline{i}(x,Ay) + i(y,Ax) + i \overline{i} (y,Ay) - (x,Ax) + \overline{i}(x,Ay) + i(y,Ax) - i \overline{i} (y,Ay)] = \\
= 2 (x,Ay) + 2 (y,Ax) + i[2\overline{i}(x,Ay) + 2i(y,Ax)] = 2(x,Ay)+2(y,Ax) + 2(x,Ay) - 2(y,Ax) = 4(x,Ay)$.\\
За теоремою Ріса, можна встановити функціонал $\varphi(x,y) = 4(x,Ay)$.
\end{proof}

\begin{remark}
В дійсному випадку поляризаційна тотожність має такий вигляд:\\
$S(x+y) - S(x-y) = 2 (\varphi(x,y) + \varphi(y,x))$.
\end{remark}

\begin{corollary}
В комплексному випадку в нас виникає однозначна відповідність\\
 \begin{tikzcd}[every arrow/.append style={shift left}, column sep = huge]
\underset{\textrm{лін. опер.}}{A} \arrow{r}{\textrm{побудова}} & \underset{\textrm{півторалінійний ф-л}}{\varphi(x,y) = (x,Ay)} \arrow{r}{\textrm{побудова}} \arrow{l}{\textrm{Th. Ріса}} & \underset{\textrm{квадр. форма}}{S(x) = \varphi(x,x)} \arrow{l}{\textrm{поляр. тотожн.}} \\
\end{tikzcd}
\end{corollary}

Випадок дійсно значної квадратичної форми:\\
$\forall x \in E: S(x) = \varphi(x,x) = (x,Ax) \in \mathbb{R}$\\
Покажемо, що $A$ - самоспряжений\\
$\boxed{\Rightarrow}$ Дано: $A$ - самоспряжений\\
$S(x) = (x,Ax)$\\
$\overline{S(x)} = \overline{(x,Ax)} = \overline{(Ax,x)} = (x,Ax) = S(x)$\\
Тобто $S(x) \in \mathbb{R}$
\bigskip \\
$\boxed{\Leftarrow}$ Дано: $(x,Ax) \in \mathbb{R}$\\
Скористаємось розкладом: $\forall B: E \to E: \exists ! B_1,B_2$ - самоспряжені: $B = B_1 + i B_2$ (згодом доведу)\\
$A = A_1 + i A_2$\\
$S(x) = (x,Ax)= (x,A_1x) + (x,iA_2x) = (x,A_1x) - i(x,A_2x) \in \mathbb{R}$\\
$\Rightarrow (x,A_2x) = 0 \Rightarrow A_2 = 0$\\
Отже, $A = A_1$ - самоспряжений
\bigskip \\
Заданий $E$ - евклідів простір із $(\cdot, \cdot)$ над $\mathbb{R}$\\
$S(X) = \varphi(x,x)$\\
Ми отримали рівність:\\
$S(x+y)-S(x-y) = 2\varpi(x,y) + 2\varphi(y,x)$\\
Таким чином, за квадратичною формою відновлюється симетрична частина півторалінійного (білінійного) функціоналу:\\
$\varphi(x,y)$ - білінійний\\
$\varphi(x,y) =  \underbrace{\dfrac{1}{2} \left(\varphi(x,y) + \varphi(y,x) \right)}_{= \varphi_{sim}} + \underbrace{\dfrac{1}{2} \left(\varphi(x,y) - \varphi(y,x) \right)}_{= \varphi_{cosim}}$\\
Маємо, що $\varphi_{sim}(y,x) = \varphi_{sim}(x,y)$ - симетричний\\
А $\varphi_{cosim}(y,x) = -\varphi_{cosim}(x,y)$ - кососиметричний\\
Отже, $\varphi = \varphi_{sim} + \varphi_{cosim}$\\
Пригадаємо, що $\varphi(x,y) = (x,Ay)$\\
$\Rightarrow \varphi_{sim}(x,y) = \dots = \left(x, \dfrac{A+A^*}{2}x \right)$\\
Отже, отримали, що в дійсному випадку квадратичній формі однозначно відповідає\\
\begin{tikzcd}[every arrow/.append style={shift left}, column sep = huge]
S(x) \arrow{r}{\textrm{поперед.}} & \varphi_{sim}(x,y) \arrow{l}{\textrm{побудова}} \arrow{r}{\textrm{Th. Ріса}} & \underset{самоспряжений}{A} \arrow{l}{\textrm{побудова}}
\end{tikzcd}
\bigskip \\
Розглянемо таку квадратичну форму:\\
$S_B(x) = (x,Bx) = (x,B_1x) + (x,B_2x) \boxed{=} $\\
$(x,B_2x) = (B_2*x,x) = -(B_2x,x) = -(x,B_2x) \Rightarrow (x,B_2x) = 0$\\
Тобто квадратична форма для кососиметричного оператору $\equiv 0$\\
$\boxed{=} (x,B_1x)$
\bigskip \\
\textbf{Приведення квадратичної форми до канонічного вигляду}\\
Заданий $E$ - евклідів простір із $(\cdot,\cdot)$\\
$S(x) = (x,Ax)$ - дійсно значна, і $A$ - самоспряжений\\
Тоді в $E$ існує ортонормований базис з власних векторів $A$ - $\{f_1,\dots,f_n\}$\\
$\mathbb{A}_f = \begin{pmatrix}
\lambda_1 & \dots & 0 \\
\vdots & \ddots & \vdots \\
0 & \dots & \lambda_n
\end{pmatrix}$\\
$U$ - перехід від розкладу за $\{f_1,\dots,f_n\}$ до розкладу в стандартний базис\\
Тоді\\
$A = U \mathbb{A}_{f} U^{-1} = U \mathbb{A}_{f} U^{*}$\\
Тоді\\
$S(x) = (x,Ax) = (x,U \mathbb{A}_{f} U^{*} x) = (U^*x, \mathbb{A}_{f} U^{*} x) =$\\
$U^* x = y = y_1f_1 + \dots + y_nf_n$\\
$= \displaystyle \left(\sum_{j=1}^n y_jf_j, \mathbb{A}_{f} \sum_{k=1}^n y_k f_k \right) \sum_{j=1}^n \sum_{k=1}^n y_j \overline{y_k} (f_j, \mathbb{A}_f f_k) = \sum_{j=1}^n \sum_{k=1}^n y_j \overline{y_k} \lambda_k (f_j,f_k) = \sum_{j=1}^n \lambda_j (y_j)^2$\\
Висновок: дійснозначна квадратична форма $S(x) = (x,Ax)$ має канонічний вигляд:\\
$S(x) = \displaystyle \sum_{j=1}^n \lambda_j (y_j)^2$\\
$U^*x = y = \displaystyle \sum_{j=1}^n y_j f_j$\\
\subsection{Знаковизначеність лінійного оператора, матриці}
Заданий $E$ - евклідів простір із $(\cdot,\cdot)$ і $A$ - самоспряжений оператор\\
(Або заданий $E = \mathbb{R}^n (\mathbb{C}^n)$ - евклідів простір із $(\cdot,\cdot)$ і $\mathbb{A}$ - симетрична матриця)\\
\defin{5.10.1.} Оператор $A$ (матриця $\mathbb{A}$) називається:\\
- \textbf{строго додатньо визначеною}, якщо $\forall x \neq 0: S(x) = (x,Ax) > 0$\\
- \textbf{строго від'ємно визначеною}, якщо $\forall x \neq 0: S(x) = (x,Ax) < 0$\\
- \textbf{невід'ємно визначеною}, якщо $\forall x \neq 0: S(x) = (x,Ax) \geq 0$\\
- \textbf{недодатньо визначеною}, якщо $\forall x \neq 0: S(x) = (x,Ax) \leq 0$\\
Позначення: $A > 0$ і т.д.
\bigskip \\
\th{5.10.2. Критерій знаковизначеності}\\
Заданий $E$ - евклідів простір із $(\cdot,\cdot)$ і $A$ - самоспряжений оператор\\
Тоді:\\
$A > 0 \iff$ всі власні числа $\lambda_j > 0$\\
$A < 0 \iff$ всі власні числа $\lambda_j < 0$\\
Для нестрогих нерівностей теж саме\\
\proof
$S(x) = (Ax,x) = \displaystyle \sum_{j=1}^n \lambda_j |y_j|^2$\\
$A > 0 \Rightarrow \forall x \Rightarrow \forall \vec{y} = \begin{pmatrix}
y_1 \\ \dots \\ y_n
\end{pmatrix}: \displaystyle \sum_{j=1}^n \lambda_j |y_j|^2 > 0 \Rightarrow$\\
$\vec{y} = \begin{pmatrix}
1 \\ 0 \\ \dots \\ 0
\end{pmatrix}, \lambda_1 > 0 \dots \vec{y} = \begin{pmatrix}
0 \\ 0 \\ \dots \\ 1
\end{pmatrix}, \lambda_n > 0$\\
А якщо $\lambda_1,\dots,\lambda_n > 0$, то автоматично $S(x) > 0$, тоді $A>0$\\
Решта випадків аналогічні \qed
\bigskip \\
\prp{5.10.3.} $A \leq 0 \iff (-A) \geq 0$\\
\textit{За визначенням}
\bigskip \\
\th{5.10.4. Критерій Сільвестра} (знаковизначеність матриці)\\
Задана матриця $A$ - симетрична\\
Обчислимо $\Delta_1 = M_1^1, \Delta_2 = M_{12}^{12}, \dots, \Delta_k = M^{12\dots k}_{12 \dots, k}$ - головні кутові мінори\\
$A>0 \iff \forall k: \Delta_k > 0$\\
$A < 0 \iff \forall k: (-1)^k \Delta_k > 0$\\
Для нестрогої визначеності теж саме\\
\textit{Без доведення (див. підручник Гантмахер "Теория матриц")}
\bigskip \\
\fi

\iffalse
\subsection{Зведення кривих та поверхонь 2-го порядку до канонічного вигляду}
Маємо загальний вигляд кривої другого порядку:
\begin{align*}
a_{11}x^2 + 2a_{12}xy + a_{22}y^2 + b_1x + b_2y + c = 0
\end{align*}
Розглянемо квадратну частину:\\
$a_{11}x^2 + 2a_{12}xy + a_{22}y^2 = \left( \begin{pmatrix}
 a_{11} & a_{12} \\
 a_{12} & a_{22}
\end{pmatrix} \begin{pmatrix}
x \\ y
\end{pmatrix}, \begin{pmatrix}
x \\ y
\end{pmatrix} \right)$ - квадратична форма, а матриця $\mathbb{A} = \begin{pmatrix}
a_{11} & a_{12} \\
a_{12} & a_{22}
\end{pmatrix}$ - відповідно самоспряжена матриця.
\bigskip \\
Також маємо загальний вигляд поверхну другого порядку:
\begin{align*}
a_{11}x^2 + a_{22}y^2 + a_{33}z^2 + 2a_{12}xy + 2a_{13}xz + 2a_{23}yz + b_1x + b_2y + b_3z + c = 0
\end{align*}
Аналогічно розглянемо квадратну частину:\\
$a_{11}x^2 + a_{22}y^2 + a_{33}z^2 + 2a_{12}xy + 2a_{13}xz = \left(\begin{pmatrix}
a_{11} & a_{12} & a_{13} \\
a_{12} & a_{22} & a_{23} \\
a_{13} & a_{23} & a_{33}
\end{pmatrix} \begin{pmatrix}
x \\ y \\ z
\end{pmatrix}, \begin{pmatrix}
x \\ y \\ z
\end{pmatrix} \right)$ - квадратична форма, зі самоспряженою матрицею $\mathbb{A} = \begin{pmatrix}
a_{11} & a_{12} & a_{13} \\
a_{12} & a_{22} & a_{23} \\
a_{13} & a_{23} & a_{33}
\end{pmatrix}$.
\bigskip \\
Тоді рівняння кривої (випадок $\mathbb{R}^2$) та поверхні (випадок $\mathbb{R}^3$) другого порядку записується таким чином:\\
$(\vec{x}, A \vec{x}) + (\vec{b}, \vec{x}) + c = 0$.\\
Це рівняння зведемо до канонічного вигляду:\\
1. Оскільки $A$ - самоспряжений, то існує ортонормований базис із власних векторів $A$.\\
$U$ - унітарна матриця переходу: $\mathbb{A} = U \mathbb{A}_{\textrm{diag}} U^*$. Тоді маємо:\\
$(A\vec{x}, \vec{x}) = (U A_{\textrm{diag}} U^* \vec{x}, \vec{x}) = (A_{\textrm{diag}} U^* \vec{x}, U^* \vec{x}) = (A_{\textrm{diag}}\vec{y}, \vec{y}) = \left[ \begin{gathered} \lambda_1 y_1^2 + \lambda_2 y_2^2 \\
\lambda_1 y_1^2 + \lambda_2 y_2^2 + \lambda_3 y_3^2
 \end{gathered} \right.$

2. $(\vec{b}, \vec{x}) = (U^* \vec{b}, U^* \vec{x}) = (U^*\vec{b}, \vec{y})$.
Знаходимо $\vec{\widetilde{b}} = U^* \vec{b}$, тоді\\
$(\vec{b}, \vec{x}) = (\vec{\widetilde{b}}, \vec{y}) = \left[ \begin{gathered} \widetilde{b_1}y_1 + \widetilde{b_2}y_2 \\ \widetilde{b_1}y_1 + \widetilde{b_2}y_2 + \widetilde{b_3}y_3 \end{gathered} \right.$\\
Остаточно отримаємо:\\
$\lambda_1 y_1^2 + \lambda_2 y_2^2 + \lambda_3 y_3^2 + \widetilde{b_1}y_1 + \widetilde{b_2}y_2 + \widetilde{b_3}y_3 + c = 0$.
\bigskip \\
Нехай $\lambda_1, \lambda_2, \lambda_3 \neq 0$. Тоді виділимо повні квадрати:\\
$\lambda_1 \left( y_1 + \dfrac{\widetilde{b_1}}{2\lambda_1} \right)^2 + \lambda_2 \left( y_2 + \dfrac{\widetilde{b_2}}{2\lambda_2} \right)^2  + \lambda_3 \left( y_3 + \dfrac{\widetilde{b_3}}{2\lambda_3} \right)^2 = - c + \dfrac{\widetilde{b_1}^2}{4 \lambda_1} + \dfrac{\widetilde{b_2}^2}{4 \lambda_2} + \dfrac{\widetilde{b_3}^2}{4 \lambda_3}$\\
Зробимо заміни:\\
$
z_1 = y_1 + \dfrac{\widetilde{b_1}}{2\lambda_1} \hspace{0.5cm}
z_2 = y_2 + \dfrac{\widetilde{b_2}}{2\lambda_2} \hspace{0.5cm}
z_3 = y_3 + \dfrac{\widetilde{b_3}}{2\lambda_3} \hspace{0.5cm}
\widetilde{c} = -c + \dfrac{\widetilde{b_1}^2}{4 \lambda_1} + \dfrac{\widetilde{b_2}^2}{4 \lambda_2} + \dfrac{\widetilde{b_3}^2}{4 \lambda_3}
$\\
Отримаємо:\\
$\lambda_1 z_1^2 + \lambda_2 z_2^2 + \lambda_3 z_3^2 = \widetilde{c}$.
\bigskip \\
Випадок, коли лише $\lambda_1 = 0$. Тоді робимо ті самі процедури та отримаємо:\\
$\lambda_2 \left( y_2 + \dfrac{\widetilde{b_2}}{2\lambda_2} \right)^2  + \lambda_3 \left( y_3 + \dfrac{\widetilde{b_3}}{2\lambda_3} \right)^2 + \widetilde{b_1}y_1 = - c + \dfrac{\widetilde{b_1}^2}{4 \lambda_1} + \dfrac{\widetilde{b_2}^2}{4 \lambda_2} + \dfrac{\widetilde{b_3}^2}{4 \lambda_3}$\\
Зробимо заміни:\\
$
z_2 = y_2 + \dfrac{\widetilde{b_2}}{2\lambda_2} \hspace{0.5cm}
z_3 = y_3 + \dfrac{\widetilde{b_3}}{2\lambda_3} \hspace{0.5cm}
\widetilde{c} = -c + \dfrac{\widetilde{b_2}^2}{4 \lambda_2} + \dfrac{\widetilde{b_3}^2}{4 \lambda_3} \hspace{0.5cm}
z_1 = y_1 - \dfrac{\widetilde{c}}{\widetilde{b_1}}
$\\
$\lambda_2 z_2^2 + \lambda_3 z_3^2 + \widetilde{b_1}z_1 = 0$ - а це є параболоїдом.\\
Випадок $\widetilde{b_1} = 0$: в $\mathbb{R}^3$ - циліндр, а в $\mathbb{R}^2$ - пара прямих.
\bigskip \\
Випадок, коли $\lambda_1 = \lambda_2 = 0$ (це вже лише для $\mathbb{R}^3$):\\
$\lambda_3 y_3^2 + \widetilde{b_1}y_1 + \widetilde{b_2}y_2 + \widetilde{b_3}y_3 + c = 0$\\
Виникає одна мрія: $\widetilde{b_1} = 0$\\
$\vec{\widetilde{b}} = U^* \vec{b}$\\
$U = (\vec{f_1}, \vec{f_2}, \vec{f_3})$, де $\vec{f_1}, \vec{f_2}$ - власні числа для $\lambda_1 = \lambda_2$, а $\vec{f_3}$ - для $\lambda_3 \neq 0$\\
$U^* \vec{b} = \begin{pmatrix}
 \overset{\leftarrow}{f_1} \cdot \vec{b} \\ \overset{\leftarrow}{f_2} \cdot \vec{b} \\ \overset{\leftarrow}{f_3} \cdot \vec{b}
\end{pmatrix} = \begin{pmatrix}
 (\vec{f_1}, \vec{b}) \\ (\vec{f_2}, \vec{b}) \\ (\vec{f_3}, \vec{b})
\end{pmatrix} = \begin{pmatrix}
 \widetilde{b_1} \\ \widetilde{b_2} \\ \widetilde{b_3}
\end{pmatrix}$\\
Щоб здійснити мрію, треба, щоб $(\vec{f_1}, \vec{b}) = 0$\\
Знаходимо $\vec{f_3}$ - власний для $\lambda_3$\\
Тоді $\linspan\{e_3\}^{\perp}$ - простір власних векторів для $\lambda_1 = \lambda_2 = 0$\\
Тоді $\vec{f_1} = [\vec{f_3}, \vec{b}] \perp \vec{b} \Rightarrow \widetilde{b_1} = 0$\\
$f_1$ - власний для $\lambda = 0$ та $f_2 = [\vec{f_1}, \vec{f_3}] \perp \vec{f_1}, \vec{f_3}$\\
Параметризуємо $\vec{f_1}, \vec{f_2}, \vec{f_3}$ - отримаємо ортонормований базис\\
$U$ - матриця перехода $\vec{\widetilde{b}} = \begin{pmatrix}
0 \\ \widetilde{b_2} \\ \widetilde{b_3}
\end{pmatrix}$\\
Тоді:\\
$\lambda_3 y_3^2 + \widetilde{b_3} y_3 + \widetilde{b_2} y_2 + c = 0$\\
$\lambda_3 \left(y_3 + \dfrac{\widetilde{b_3}}{2 \lambda_3} \right)^2 + \widetilde{b_2} \left(y_2 + \dfrac{c - \dfrac{\widetilde{b_3}^2}{4\lambda_3}}{\widetilde{b_2}} \right) = 0$\\
Зробимо заміни:\\
$z_1 = y_1 \hspace{0.5cm}
z_2 = y_2 + \dfrac{c - \dfrac{\widetilde{b_3}^2}{4\lambda_3}}{\widetilde{b_2}} \hspace{0.5cm}
z_3 = y_3 + \dfrac{\widetilde{b_3}}{2 \lambda_3}
$\\
Отримаємо:\\
$\lambda_3 z_3^2 + \widetilde{b_2} z_2 = 0$ - а це є параболічним циліндром\\
При $\widetilde{b_2} = 0$ отримаємо пару площин.\\
$\scaleobj{5}{\blacksquare}$
\fi

\iffalse
\section{Додаткові факти}
Дани розділ присвячений виключно для тих, хто володіє фактами з теорії метричних просторів та функціонального аналізу. Всі нові означення, які були наведені там, я нагадувати не буду. Із ними можна ознайомитися в іншому pdf файлі.
\bigskip \\
Перша теорема дає пояснення, чому в лінійній алгебрі (де досліджуються скінченні простори) не вводилися норми, на відміну від функціонального аналізу (де вже нескінченні простори).
\begin{theorem}
У кожному скінченно вимірному просторі всі норми еквівалентні.
\end{theorem}

\begin{proof}
Ми знаємо, що всі скінченні простори ізоморфні до $\mathbb{R}^d$, тому будемо тут досліджувати це питання. Досить буде довести, що всі норми еквівалентні до $\| \cdot \|_2$.\\
Нехай $\{\vec{e}_1,\dots,\vec{e}_d\}$ -- стандартний базис $\mathbb{R}^d$, тоді звідси $\vec{x} = \displaystyle\sum_{i=1}^d x_i \vec{e}_i$. Оберемо довільну норму $\|\cdot \|$.\\
$\displaystyle\left\| \sum_{i=1}^d x_i \vec{e}_i
 \right\| \leq \sum_{i=1}^d \| x_i e_i \| = \sum_{i=1}^d |x_i| \|e_i\| = \sqrt{\left( \sum_{i=1}^d |x_i| \|\vec{e}_i\| \right)^2} \overset{\text{К-Б}}{\leq} \sqrt{\sum_{i=1}^d \|e_i\|^2} \sqrt{\sum_{j=1}^d |x_j|^2} \leq \sqrt{\sum_{i=1}^d \|e_i\|^2} \sqrt{\sum_{j=1}^\infty |x_j|^2} = \sqrt{\sum_{i=1}^d \|e_i\|^2} \|\vec{x}\|_2 = M \|\vec{x}\|_2$.\\
Зауважимо, що $M \in \mathbb{R}_{\geq 0}$ та не залежить від $\vec{x}$. Отже, $\|\vec{x}\| \leq M \|\vec{x}\|_2$.
\bigskip \\
Розглянемо тепер $\mathcal{S}$ -- одинична сфера на $(\mathbb{R}^d, \|\cdot \|_2)$. Відомо, що $S$ -- замкнена множина та обмежена. Тож за лемою Гайне-Бореля, $S$ -- компактна множина. Відомо, що відображення $\| \cdot \| \colon \mathcal{S} \to \mathbb{R}_{\geq 0}$ -- неперервне відображення, тож вона досягає найменшого значення $m$ для деякого $\vec{y} \in S$.\\
Припустимо $m = 0$, тоді звідси $\|\vec{y}\| = 0 \implies \vec{y} = \vec{0} \implies \vec{y} \notin \mathcal{S}$ -- неможливо. Отже, $m > 0$.\\
Значить, $\forall \vec{y} \in \mathbb{R}^d: \|\vec{y}\|_2 = 1: \|y\| \geq m$. Треба довести те саме для інших векторів.\\
Якщо $\vec{x} = \vec{0}$, то це виконано. Тому $\vec{x} \neq \vec{0}$. Покладемо вектор $\vec{y} = \dfrac{\vec{x}}{\| \vec{x} \|_2}$, причому $\|\vec{y}\|_2 = 1$. Із цього випливає, що $\| \vec{y} \|_2 \leq m \implies m \|\vec{x}\|_2 \leq \|\vec{x}\|$.
\bigskip \\
Всі інші норми будуть еквівалентними в силу транзитивності.
\end{proof}

\begin{theorem}
Лінійні оператори між скінченно вимірними векторними просторами -- обмежені.\\
\textit{Це, насправді, наслідок твердження нижче.}
\end{theorem}

\begin{proposition}
Задано $X,Y$ -- нормовані простори, $\dim X < \infty$ та $A \in \mathcal{L}(X,Y)$. Тоді $A \in \mathcal{B}(X,Y)$.
\end{proposition}

\begin{proof}
Дійсно, нехай $\{e_1,\dots,e_n\}$ -- базис $X$, нехай на неї стоїть норма $\|x\|_2$, тоді маємо наступне:\\
$\|Ax\| = \|A (x_1 e_1 + \dots + x_n e_n) \| = \| x_1 A e_1 + \dots + x_n A e_n \| \leq |x_1| \|A e_1 \| + \dots + |x_n| \|A e_n \| \leq \\
\leq \sqrt{|x_1|^2 + \dots + |x_n|^2} \sqrt{\|Ae_1\|^2 + \dots + \|Ae_n\|^2} = C \|x\|_2$.\\
Якби була би інша норма $\| \cdot \|$, то вона еквівалентна $\| \cdot \|_2$, а тому обмеженість зберігається.
\end{proof}

\begin{theorem}
Задано $E$ -- нормований простір.\\
$\dim E < \infty \iff$ кожна підмножина $E$ -- передкомпактна.
\end{theorem}

\begin{proof}
\leftproof Дано: $\dim E = \infty$. Ми хочемо довести, що замкнена куля $B[0;1]$ не буде компактом.\\
У якості $x_1$ оберемо будь-який вектор одиничної сфери та покладемо $G_1 = \linspan\{x_1\}$. За лемою Ріса, знайдеться вектор $x_2 = y_{\frac{1}{2}} \notin G_1$ такий, що $\|x_2\| = 1$ та $\forall x \in G_1: \|x_2-x\| > \dfrac{1}{2}$, зокрема звідси $\|x_2 - x_1\| > \dfrac{1}{2}$.\\
Покладемо $G_2 = \linspan\{x_1,x_2\}$. Тоді знову за лемою Ріса, знайдеться вектор $x_3 = y_{\frac{1}{2}}$ такий, що $\|x_3\| = 1$ та $\forall x \in G_2: \|x_3-x\| > \dfrac{1}{2}$. Зокрема $\|x_3 - x_1\| > \dfrac{1}{2},\ \|x_3 - x_2\| > \dfrac{1}{2}$.\\
\vdots \\
У силу того, що $\dim E = \infty$, то цей процес буде продовжуватися. Побудуємо послідовність $(x_n)_{n=1}^\infty \subset B[1;0]$ таку, що $\forall m,n \geq 1: \| x_m - x_n \| > \dfrac{1}{2}$. Така обмежена послідовність не містить збіжної підпослідовності.
\bigskip \\
\rightproof Дано: $\dim E < \infty$, тоді відомо, що $E \cong \mathbb{C}^n$. Будь-яка обмежена підмножина $\mathbb{R}^n$ -- передкомпактна за Больцано-Ваєрштраса.
\end{proof}
\newpage
\fi

\iffalse
\section*{Постскриптум}
\subsection*{Про теорему Ріса}
Під час розгляду евклідових чи унітарних просторів я навів теорему в такій формі:
\begin{theorem}[Теорема Ріса]
Задано $(E,(\cdot,\cdot))$ -- унітарний (евклідів) простір та лінійний функціонал $\varphi \colon E \to \mathbb{C} (\mathbb{R})$. Тоді $\exists ! f \in E: \forall x \in E: \varphi(x) = (x,f)$.
\end{theorem}
\noindent
Чесна та правильна версія виглядає саме таким чином:
\begin{theorem}[Теорема Ріса]
Задано $(H,(\cdot,\cdot))$ -- гілбертів простір та лінійний обмежений функціонал $l \colon H \to \mathbb{C} (\mathbb{R})$. Тоді $\exists ! f \in H: \forall x \in H: l(x) = (x,f)$.
\end{theorem}
\noindent
По-перше, тут тепер вимагається, щоб простір був гілбертовим. По суті, різниця між гілбертовим та унітарним полягає лише в тому, що унітарний обов'зяково скінченновимірний простір, водночас в гілбертовому це не обов'язково.\\
По-друге, крім лінійності функціонала тепер ще накладається умова обмеженості. Будь-який функціонал, що визначений на скінченномовимірному просторі, автоматично обмежений. Тому я цю вимогу окремо не розписував.\\
По-третє, важливе зауваження: існуючий єдиний $f$ не залежить більше ні від чого, крім від вибору функціоналу (коли ми доводили цю теорему в цьому пдф, то там залежність ще була від базису). Це означає, що $H \cong H^*$ канонічним чином та прийнято писати тоді $H = H^*$.
\fi

\iffalse
\section*{Використані джерела}
\begin{enumerate}
\item \href{https://www.youtube.com/watch?v=Lcy5VDYMM4Q&list=PLgAugiET8rrIGTphP-IfbxixpAZDO3UHC&index=3}{DanielChanMaths: Properties of tensors products of vector spaces}
\item \href{https://www.youtube.com/watch?v=KnSZBjnd_74&ab_channel=MuPrimeMath}{Mu Prime Math: A Concrete Introduction to Tensor Products}
\item \href{https://www.youtube.com/watch?v=vZzZhdLC_YQ&t=0s&ab_channel=MuPrimeMath}{Mu Prime Math: Complete Derivation: Universal Property of the Tensor Product}
\item \href{https://www.youtube.com/watch?v=5i8tp0rD2EA&ab_channel=MuPrimeMath}{My Prime Math: Tensor Product Basis With the Universal Property}
\item \href{https://www.youtube.com/watch?v=K7f2pCQ3p3U}{Michael Penn: What is a tensor anyway?? (from a mathematician)}
\end{enumerate}
\fi